\section{2022 Analysis \& Differential Equations}

\subsection{Exam}

\begin{exercise}
\label{analysis:2022:1}
For $n\geq 1$, we consider the integral
\[
I_n = \int_{[0,1]^n} \frac{n}{\frac{1}{x_1}+\frac{1}{x_2}+\cdots+\frac{1}{x_n}}\,dx_1\dots dx_n.
\]
Prove that $\lim_{n\to\infty} I_n$ exists.
\end{exercise}

\begin{solution}
\textbf{Prerequisites:}
To understand this problem, you should be familiar with the concept of an expectation in probability (which corresponds to an integral over a probability space), the properties of the Harmonic Mean, the Strong Law of Large Numbers (SLLN) for variables with infinite mean, and the Dominated Convergence Theorem (DCT).


\textbf{Steps:}
\begin{enumerate}
    \item \textbf{Probabilistic Interpretation:} We recognize the integral $I_n$ as the expected value of the harmonic mean of $n$ independent and identically distributed (i.i.d.) random variables $X_1, X_2, \dots, X_n$, each uniformly distributed on the interval $[0,1]$. Thus, $I_n = E[H_n]$, where $H_n = n / \sum_{i=1}^n (1/X_i)$.

    \item \textbf{Analysis of the Reciprocals:} Let $Y_i = 1/X_i$. We calculate the expected value of these transformed variables: $E[Y_1] = \int_0^1 (1/x) dx$, which is infinite. This means that, on average, the reciprocals of our numbers are extremely large.

    \item \textbf{Applying the Law of Large Numbers:} The Strong Law of Large Numbers tells us that the average of these reciprocals, $\frac{1}{n} \sum_{i=1}^n Y_i$, will almost surely approach the mean $E[Y_1] = \infty$ as $n$ grows. Consequently, its reciprocal—the harmonic mean $H_n$—must almost surely approach $1/\infty = 0$.

    \item \textbf{Passing to the Limit:} We observe that since $X_i \in [0,1]$, each $1/X_i \geq 1$, which implies the sum $\sum 1/X_i \geq n$. Therefore, the harmonic mean $H_n$ is always between $0$ and $1$. Since $H_n$ is bounded by the integrable constant function $1$, we can use the Dominated Convergence Theorem to conclude that $\lim_{n\to\infty} E[H_n] = E[\lim_{n\to\infty} H_n] = E[0] = 0$.
\end{enumerate}


\textbf{Intuition:}
The harmonic mean is very sensitive to small numbers in a dataset. Because we are picking numbers from $[0,1]$, as we pick more and more numbers, we are almost guaranteed to pick some that are extremely close to zero. These tiny numbers have massive reciprocals, which dominate the sum in the denominator and pull the entire harmonic mean down to zero.
\end{solution}

\begin{exercise}
\label{analysis:2022:2}
Let $U\subset\mathbb{C}$ be a non-empty open set and $f:U\to U$ be a non-constant holomorphic function. Prove that, if $f\circ f = f$, then $f(z)\equiv z$ for all $z\in U$.
\end{exercise}

\begin{solution}
\textbf{Prerequisites:}
Knowledge of basic complex analysis is required, specifically the Open Mapping Theorem (non-constant holomorphic functions map open sets to open sets) and the Identity Theorem (if two holomorphic functions agree on an open set or a set with a limit point, they agree everywhere on a connected domain).


\textbf{Steps:}
\begin{enumerate}
    \item \textbf{Statement risk and connectedness convention:} The conclusion is correct when $U$ is a domain, i.e. a connected open set. This connectedness assumption is necessary. If $U$ is disconnected, take two disjoint disks $U_1,U_2$, define $f(z)=z$ on $U_1$, and define $f$ to be a constant point of $U_1$ on $U_2$. Then $f:U_1\cup U_2\to U_1\cup U_2$ is holomorphic, non-constant, and satisfies $f\circ f=f$, but it is not the identity on $U_2$.

    \item \textbf{Image Analysis:} Let $V = f(U)$ be the image of $U$. By the Open Mapping Theorem, because $f$ is non-constant and holomorphic, $V$ must be an open subset of $U$.

    \item \textbf{Fixed Point Property:} The condition $f(f(z)) = f(z)$ tells us that for any point $w$ in the image $V$, $f(w) = w$. In other words, $f$ acts as the identity function on the set $V$.

    \item \textbf{Applying the Identity Theorem:} Define a new function $g(z) = f(z) - z$. We have shown that $g(z) = 0$ for all $z$ in the non-empty open set $V$. If $U$ is connected, the Identity Theorem forces $g(z)=0$ on all of $U$.

    \item \textbf{Conclusion:} Therefore, under the domain assumption, $f(z)=z$ for all $z\in U$.
\end{enumerate}


\textbf{Intuition:}
Imagine $f$ as a way of moving points around in the set $U$. The rule $f(f(z)) = f(z)$ says that after one application, the point lands in a region where applying $f$ again changes nothing. On a connected domain, holomorphic rigidity propagates this open fixed region to the whole domain.
\end{solution}

\begin{exercise}
\label{analysis:2022:3}
Let $X\subset\mathbb{R}$ be a set with positive (Lebesgue) measure. Show that we can find an arithmetic progression of $2022$ terms in $X$, i.e., there exists $x_1,\dots,x_{2022}\in X$ so that the $x_{i+1}-x_i$'s are all equal and positive, $i=1,\ldots,2021$.
\end{exercise}

\begin{solution}
\textbf{Prerequisites:}
You should understand Lebesgue measure, indicator functions ($1_X$), and the continuity of the translation operator in $L^p$ spaces (specifically, how $\int |f(x+d) - f(x)|^p dx \to 0$ as $d \to 0$).


\textbf{Steps:}
\begin{enumerate}
    \item \textbf{Reduce to finite measure:} Since $X$ has positive measure, choose $R>0$ such that
    \[
    E:=X\cap[-R,R]
    \]
    has finite positive measure.

    \item \textbf{Look for a common overlap:} We want
    \[
    x,\ x+d,\ \ldots,\ x+2021d
    \]
    to all lie in $E$. This is equivalent to
    \[
    x\in \bigcap_{j=0}^{2021}(E-jd).
    \]

    \item \textbf{Use translation continuity:} Since $1_E\in L^1(\mathbb{R})$, translations are continuous in $L^1$:
    \[
    m(E\triangle(E-jd))
    =
    \|1_E-1_{E-jd}\|_{L^1}
    \to0
    \]
    as $d\to0$, for each fixed $j$.

    \item \textbf{Keep the intersection positive:} By the union bound,
    \[
    m\left(E\setminus\bigcap_{j=0}^{2021}(E-jd)\right)
    \leq
    \sum_{j=0}^{2021}m(E\setminus(E-jd))
    \leq
    \sum_{j=0}^{2021}m(E\triangle(E-jd)).
    \]
    For all sufficiently small positive $d$, the right side is less than $m(E)$. Hence the intersection has positive measure, so there exists an $x$ in it. Then $x+jd\in E\subset X$ for $j=0,\ldots,2021$, giving the required arithmetic progression.
\end{enumerate}


\textbf{Intuition:}
Think of $X$ as a collection of occupied spots on the number line. A finite positive-measure piece of $X$ still overlaps strongly with several tiny translates of itself, and any point in that common overlap starts the desired arithmetic progression.
\end{solution}

\begin{exercise}
\label{analysis:2022:4}
Let $C([0,1])$ be the space of all continuous $\mathbb{R}$-valued functions equipped with $L^{\infty}$-norm. Let $\mathbf{P}\subset C([0,1])$ be a closed linear subspace. Assume that the elements of $\mathbf{P}$ are polynomials. Prove that $\dim\mathbf{P}<\infty$.
\end{exercise}

\begin{solution}
\textbf{Prerequisites:}
To solve this problem, you need to be familiar with Banach spaces (complete normed vector spaces), the Baire Category Theorem (which discusses the structure of complete metric spaces as unions of sets), and the fact that finite-dimensional subspaces of a normed space are always closed.


\textbf{Steps:}
\begin{enumerate}
    \item \textbf{Stratify by Degree:} Let $\mathcal{P}_n$ be the space of all polynomials of degree at most $n$. These are finite-dimensional spaces. We define $\mathbf{P}_n = \mathbf{P} \cap \mathcal{P}_n$. Since $\mathcal{P}_n$ is closed and $\mathbf{P}$ is closed, their intersection $\mathbf{P}_n$ is also a closed subspace.

    \item \textbf{Countable Union:} By the assumption that every element of $\mathbf{P}$ is a polynomial, we can write $\mathbf{P} = \bigcup_{n=0}^\infty \mathbf{P}_n$. This means the space $\mathbf{P}$ is built out of a countable collection of closed sets.

    \item \textbf{Applying Baire Category Theorem:} Since $\mathbf{P}$ is a closed subspace of the complete space $C([0,1])$, $\mathbf{P}$ itself is a complete metric space. The Baire Category Theorem states that if a complete space is a countable union of closed sets, at least one of those sets (say $\mathbf{P}_N$) must have a non-empty interior.

    \item \textbf{Geometric Conclusion:} In a vector space, if a linear subspace has a non-empty interior (meaning it contains an open ball), it must actually be the entire space. Therefore, $\mathbf{P} = \mathbf{P}_N$. Since $\mathbf{P}_N$ is a subspace of the finite-dimensional space $\mathcal{P}_N$, its dimension must be finite (specifically, $\dim \mathbf{P} \leq N+1$).
\end{enumerate}


\textbf{Intuition:}
Imagine the space of all polynomials as a building with infinitely many floors, where each floor $n$ contains only polynomials of degree $n$ or less. The problem says our space $\mathbf{P}$ is "trapped" inside this building and is also "complete" (closed). The Baire Category Theorem essentially says you can't fill up a solid, complete space using only a bunch of "flat" (empty interior) pieces unless at least one of those pieces is actually not flat at all. If one floor of our building already contains a whole room of our space, the linear nature of the space forces the entire space to fit within that floor.
\end{solution}

\begin{exercise}
\label{analysis:2022:5}
Let $\Omega\subset\mathbb{R}^3$ be a bounded domain with smooth boundary. Assume that $u\in C(\overline{\mathbb{R}^3\setminus\Omega})$ is a harmonic function on $\mathbb{R}^3\setminus\Omega$ so that $u|_{\partial\Omega}=1$ and $\lim_{|x|\to\infty}|u(x)|=0$.

Prove that for such $u$, $\lim_{|x|\to\infty}|x|u(x)$ exists.
\end{exercise}

\begin{solution}
\textbf{Prerequisites:}
You should know about harmonic functions, spherical harmonics, and the concept of multipole expansions (often used in physics to describe gravitational or electrostatic potentials). Familiarity with Green's identities is also helpful.


\textbf{Steps:}
\begin{enumerate}
    \item \textbf{Work outside a large ball:} Choose $R>0$ such that $\overline{\Omega}\subset B_R(0)$. On the exterior region $|x|>R$, the function $u$ is harmonic and tends to $0$ at infinity.

    \item \textbf{Use the Kelvin transform:} Define, for $0<|y|<1/R$,
    \[
    v(y):=\frac{1}{|y|}\,
    u\left(\frac{y}{|y|^2}\right).
    \]
    In dimension $3$, the Kelvin transform preserves harmonicity, so $v$ is harmonic on the punctured ball $0<|y|<1/R$.

    \item \textbf{Remove the possible singular part:} By the maximum principle on exterior annuli, the boundary condition $u=1$ and the limit $u\to0$ imply $0\leq u\leq1$ on the unbounded exterior component. Thus $v\geq0$. Bôcher's theorem for a positive harmonic function on a punctured ball gives
    \[
    v(y)=\frac{a}{|y|}+h(y),
    \]
    where $a\geq0$ and $h$ is harmonic near $0$.

    \item \textbf{Use the decay at infinity:} If $x=y/|y|^2$, then $|x|=1/|y|$ and
    \[
    u(x)=|y|v(y)=a+|y|h(y).
    \]
    Since $u(x)\to0$ as $|x|\to\infty$, we must have $a=0$. Therefore $v=h$ extends harmonically to $0$.

    \item \textbf{Take the desired limit:} The desired product is exactly
    \[
    |x|u(x)=v\left(\frac{x}{|x|^2}\right).
    \]
    As $|x|\to\infty$, the argument $x/|x|^2$ tends to $0$, so
    \[
    \lim_{|x|\to\infty}|x|u(x)=v(0)
    \]
    exists.
\end{enumerate}


\textbf{Intuition:}
Think of $u(x)$ as the voltage created by a charged conductor. The Kelvin transform turns the far-field question into a removable-singularity question at the origin; once the singular coefficient is ruled out by the decay $u(x)\to0$, the transformed function has an ordinary limit.
\end{solution}

\begin{exercise}
\label{analysis:2022:6}
Let $f(x,y)\in C^1(\mathbb{R}^2)$. We assume that there exists $C>0$ so that for all $(x,y)\in\mathbb{R}^2$, $|\frac{\partial f}{\partial y}(x,y)|\leq C$. Prove that the following ODE has a globally defined solution for all $y(0)=y_0\in\mathbb{R}$:
\[
\begin{cases}
\frac{d}{dx}y(x) = f(x,y(x)),\\
y(0) = y_0.
\end{cases}
\]

In addition, we assume that $f$ is $1$-periodic in $x$, i.e., for all $(x,y)\in\mathbb{R}^2$, we have $f(x+1,y)=f(x,y)$. Prove that if the above admits a globally defined bounded solution, then it admits a periodic solution.
\end{exercise}

\begin{solution}
\textbf{Prerequisites:}
Basic knowledge of Ordinary Differential Equations (ODEs), including the Picard-Lindelöf theorem for existence/uniqueness, the concept of "blow-up" in finite time, and the Poincaré map for periodic systems.


\textbf{Steps:}
\textbf{Part 1: Global Existence}
\begin{enumerate}
    \item \textbf{Check Lipschitz Condition:} The given bound $|\partial f / \partial y| \leq C$ means the function $f$ is "globally Lipschitz" in $y$. This is a very strong condition that ensures solutions are unique and don't behave too wildly.

    \item \textbf{Linear Growth on Finite Time Windows:} For any fixed $T>0$, continuity gives
    \[
    A_T:=\sup_{|x|\leq T}|f(x,0)|<\infty.
    \]
    The mean value theorem in the $y$ variable gives
    \[
    |f(x,y)|\leq A_T+C|y|,
    \qquad |x|\leq T.
    \]

    \item \textbf{Preventing Blow-up:} On any interval where $|x|\leq T$, Gronwall's inequality bounds $|y(x)|$ in terms of $T$, $A_T$, $C$, and $|y_0|$. Therefore no finite-time blow-up can occur inside $[-T,T]$. Since $T$ is arbitrary, the solution exists for all $x\in\mathbb{R}$.
\end{enumerate}

\textbf{Part 2: Periodic Solution}
\begin{enumerate}
    \item \textbf{The Poincaré Map:} Define a map $P$ that takes a starting value $y(0)$ and tells you where the solution ends up after one period, $y(1)$. Because the equation itself is periodic, if we can find a starting point where $y(0) = y(1)$, the solution will repeat itself forever.

    \item \textbf{Continuity and Monotonicity:} Continuous dependence on initial data makes $P$ continuous. Uniqueness prevents two solution graphs from crossing, so $P$ is increasing: if $a<b$, then $P(a)\leq P(b)$.

    \item \textbf{Use a Bounded Orbit:} Let $y_b$ be a globally defined bounded solution and set $a_n=y_b(n)$. Then $a_{n+1}=P(a_n)$ and $\{a_n\}_{n\geq0}$ is bounded. If $P(a_0)>a_0$, monotonicity implies
    \[
    a_0\leq a_1\leq a_2\leq\cdots.
    \]
    If $P(a_0)<a_0$, the sequence is monotone decreasing. If $P(a_0)=a_0$, we already have a fixed point. In all cases, boundedness gives a limit $a_*$, and continuity of $P$ gives
    \[
    P(a_*)=\lim_{n\to\infty}P(a_n)
    =\lim_{n\to\infty}a_{n+1}
    =a_*.
    \]

    \item \textbf{Conclusion:} Let $y_*(x)$ be the solution with $y_*(0)=a_*$. Since $P(a_*)=a_*$, we have $y_*(1)=y_*(0)$. The shifted function $x\mapsto y_*(x+1)$ solves the same ODE and has the same initial value at $x=0$ as $y_*$. By uniqueness,
    \[
    y_*(x+1)=y_*(x)
    \]
    for all $x$, so $y_*$ is periodic.
\end{enumerate}

\textbf{Intuition:}
For the first part, the bounded $y$-derivative prevents finite-time blow-up on every finite $x$-interval. For the second part, the Poincaré map packages one full period of the equation into a one-dimensional increasing map; a bounded orbit of such a map produces a fixed point.
\end{solution}
